\documentclass[11pt]{article}
\usepackage[a4paper,margin=20mm]{geometry}
\usepackage{amsmath,amssymb,xcolor,graphicx}
\usepackage[hypcap=false]{caption}
\usepackage[hidelinks,bookmarksnumbered,bookmarksopen]{hyperref}
\hypersetup{pdftitle={Thermodynamics 3 — The laws and entropy},pdfauthor={Lecturer: Dr S. Lindqvist},
  pdfcreator={KherveNote}}
\renewcommand\labelitemii{\textbullet}
\renewcommand\labelitemiii{\textbullet}
\renewcommand\labelitemiv{\textbullet}
\renewcommand\labelenumii{\arabic{enumi}.\arabic{enumii}.}
\renewcommand\labelenumiii{\arabic{enumi}.\arabic{enumii}.\arabic{enumiii}.}
\renewcommand\labelenumiv{\arabic{enumi}.\arabic{enumii}.\arabic{enumiii}.\arabic{enumiv}.}
\setlength{\parindent}{0pt}
\setlength{\parskip}{0.6em}
\definecolor{knoteaccent}{HTML}{1A6DD8}
\definecolor{knotemuted}{HTML}{6B6B6B}
\definecolor{knotekey}{HTML}{D96B00}
\newcommand\knotetime[1]{\leavevmode\llap{\color{knotemuted}\scriptsize #1\hspace{1em}}}
\newcommand\knotekey[2][]{\par\noindent#1\colorbox{knotekey!12}{\parbox{\dimexpr\linewidth-2\fboxsep}{%
  \textbf{\color{knotekey}$\star$ Key point.} #2}}\par}
\newcommand\knotequestion[2][]{\par\noindent#1{\color{knoteaccent}\textbf{?}}~\emph{#2}\par}
\newenvironment{knotetranscript}{\par\color{knotemuted}\small}{\par}
\newcommand\knotesummary[1]{\par\noindent\fcolorbox{knoteaccent}{knoteaccent!6}{%
  \parbox{\dimexpr\linewidth-2\fboxsep-2\fboxrule}{\textbf{Summary}\par #1}}\par}
\newenvironment{knotechunk}{}{}
\pagestyle{plain}
\begin{document}
\begin{knotechunk}
{\color{knoteaccent}\rule{\linewidth}{1.5pt}}\par
{\LARGE\bfseries Thermodynamics 3 — The laws and entropy}\par
{\large Lecturer: Dr S. Lindqvist}\hfill{\color{knotemuted}07 October 2026 \textperiodcentered{} Chemistry building, LT2}\par
{\color{knoteaccent}\rule{\linewidth}{0.6pt}}\par

\section{The first law}

\knotetime{01:25}Energy is conserved: $\Delta U = q + w$ (heat and work done *on* the system).

\begin{itemize}\setlength{\itemsep}{0pt}\setlength{\parskip}{0pt}
  \item at constant volume $\Delta U = q_V$; at constant pressure $\Delta H = q_p$, with $H = U + pV$
  \item state functions ($U$, $H$, $S$, $G$) depend only on the state, not the path; $q$ and $w$ do not
\end{itemize}
\end{knotechunk}

\begin{knotechunk}
\section{Entropy and the second law}

\knotetime{10:25}\[dS = \frac{\delta q_{\mathrm{rev}}}{T}, \qquad \Delta S_{\mathrm{universe}} \geq 0\] Statistically, $S = k_B \ln W$ — the number of microstates.

\subsection{Worked example: melting ice}

\knotetime{11:15}$\Delta H_{\mathrm{fus}} = 6.01$ kJ mol$^{-1}$ at 273.15 K, so $\Delta S_{\mathrm{fus}} = 6010 / 273.15 \approx 22.0$ J K$^{-1}$ mol$^{-1}$.

\knotekey[\knotetime{11:40}]{Heat engines: no engine between $T_h$ and $T_c$ beats the Carnot efficiency $\eta = 1 - T_c/T_h$.}
\end{knotechunk}

\begin{knotechunk}
\section{Gibbs energy}

\knotetime{24:25}At constant $T$ and $p$, the second law becomes a condition on the system alone: \[G = H - TS, \qquad \Delta G = \Delta H - T\Delta S < 0 \text{ for a spontaneous change}\]

\begin{enumerate}\setlength{\itemsep}{0pt}\setlength{\parskip}{0pt}
  \item $\Delta H < 0$, $\Delta S > 0$: spontaneous at all $T$
  \item $\Delta H > 0$, $\Delta S > 0$: spontaneous above $T = \Delta H / \Delta S$
  \item $\Delta H < 0$, $\Delta S < 0$: spontaneous below $T = \Delta H / \Delta S$
  \item $\Delta H > 0$, $\Delta S < 0$: never
\end{enumerate}
\begin{itemize}\setlength{\itemsep}{0pt}\setlength{\parskip}{0pt}
  \item equilibrium: $\Delta G^\circ = -RT\ln K$
\end{itemize}
\end{knotechunk}

\begin{knotechunk}
\section{The third law}

\knotetime{36:25}The entropy of a perfect crystal tends to zero as $T \to 0$ K, which gives absolute entropies.

\knotequestion[\knotetime{36:50}]{How does residual entropy (e.g. in ice or CO) fit with the third law?}
\end{knotechunk}

\begin{knotechunk}
\section{What was said}
\begin{knotetranscript}
\knotetime{10:00:30}Three laws today, and by the end you should be able to predict whether a reaction can go.

\knotetime{10:05:40}Remember the sign convention, w is work done on the system, so compression is positive.

\knotetime{10:09:00}Entropy: the classical definition is heat exchanged reversibly divided by temperature.

\knotetime{10:13:30}Boltzmann's view is more intuitive, entropy counts the number of ways the system can be arranged.

\knotetime{10:18:10}For melting ice, six point zero one kilojoules divided by two seven three gives about twenty-two joules per kelvin per mole.

\knotetime{10:22:00}Now combine both into one quantity for the system at constant temperature and pressure, that's the Gibbs energy.

\knotetime{10:27:30}When enthalpy and entropy have the same sign, temperature decides, and the crossover is delta H over delta S.

\knotetime{10:34:40}Finally the third law, which sets the zero of entropy.
\end{knotetranscript}
\end{knotechunk}
\end{document}
